\subsection{DET-09: the actual examined RH endpoint is scalar}
An adaptive source trace goes beyond searching for the name $W$. The exported theorem in \path{Supplements/Kernels/MainTheorems.lean}, declarations \texttt{allChi} and \texttt{rhOfAll}, uses the character nonvanishing theorem and \texttt{Formalism.VW.rh\_of\_every}. Following the immediate calls reaches \path{Development/AutomorphicArithmetic/CharacterBoxes/ZeroFreeAllOfTopG.lean}, \texttt{zeroFreeAll\_of\_topG}, lines 26--34, which supplies a scalar bound and \texttt{formB} to \texttt{zeroFreeOn\_of\_tiedWin\_cyc}. The precise path and all intermediate declarations are recorded in the evidence note. The mathematical types at the endpoint are explicit.

The source's smoothed von Mangoldt function is
\[
 G_\zeta(y)=\sum_{1\le n\le\lfloor y\rfloor}\Lambda(n)(y/n-1)
 -y\log y+(1+\gamma_E)y,
\]
from \path{Formalism/StatementOC.lean}, lines 11--13. This directly matches step 5 of the note. A level $\mathcal L$ defines a finite positive real array $w_\psi$ and the scalar
\[
 G_{\mathcal L}(y)=G_\zeta(y)+\sum_{\psi\in\mathcal L.\mathrm{Nt}}w_\psi\Re g_{\psi.\chi}(y).
\]
For nontrivial primitive $\chi$, the function $g_\chi$ is the smoothed character von Mangoldt sum with the correction $yL'(1,\chi)/L(1,\chi)$. The exact definitions are \path{Formalism/SharpRestriction/RamificationFrobenius/TotientGaloisClasses/Analytic.lean}, lines 9--24, and \path{ExplicitFormulaErrorOfHostConditions.lean} in the same directory, lines 32--41. The weights equal either $1$ or $(3\mathcal L.o+1)^{-1}$ with $\mathcal L.o\in\mathbb N$. Their positivity does not assert $G_{\mathcal L}\ge0$.

\path{Development/AutomorphicArithmetic/PrimitiveCosets/CoreGm.lean}, \texttt{core\_gm}, lines 12--15 and 55--74, has a one-sided real conclusion
\[
 G_{\mathcal L}(x)\le x^{1/2+(e_D+e_L+7)\varepsilon}
\]
under its specified arithmetic data. In the family version, \path{Formalism/SharpRestriction/RamificationFrobenius/TotientGaloisClasses/TopWin.lean}, lines 23--32 and 68--93, the exponent approaches $1/2$ as $\varepsilon\downarrow0$, with bounds on power windows $[Y,Y^\lambda]$, $\lambda>1$, at arbitrarily large $Y$. The definitions of \texttt{HasWindows}, \texttt{FormBAt} and \texttt{FormB} in \path{Window.lean} in the same directory, lines 7--8 and 30--33, state the scalar analytic bridge: such $x^\theta$ bounds imply nonvanishing for $\Re s>\theta$ in the retained character family. Choosing $1/2<\theta<\min(\Re s,1)$ excludes a putative right-half-plane zero; functional-equation reflection then gives the critical line.

The implementation of \texttt{formB} is \path{Development/AutomorphicArithmetic/PrimitiveCosets/Imports.lean}, lines 8--19. It invokes \texttt{zeroFree\_of\_rpow\_window} in \path{Development/AutomorphicArithmetic/DedekindLFunctions/LandauTheorems/ZeroFreeOfRpowWindow.lean}, lines 11--14, and \texttt{lemmaL} in the adjacent \path{lemmaL.lean}, lines 13--22. The relevant \texttt{PoleData}, defined in \path{Development/AutomorphicArithmetic/DedekindLFunctions/LandauZeroFreeCriterion.lean}, lines 15--23, is data for a scalar meromorphic Laplace transform
\[
 \Phi(s)=\int_0^\infty G_F(e^u)e^{-su}\,du\qquad(\Re s>2),
\]
where $F$ is a finite weighted character family. It is not an operator trace.

\paragraph{Checked residue calculation.}
The adjacent \path{LandauTheorems/poleData.lean}, lines 137--144, 250--261 and 344--383, gives
\[
 c_\rho=\frac{S_\rho}{\rho(1-\rho)},\qquad
 S_\rho=\sum_i\frac{\alpha_i}{2}\bigl(m_i(\rho)+m'_i(\rho)\bigr)>0,
\]
where the $\alpha_i$ are positive real character weights and the multiplicities are nonnegative integers. The membership of $\rho$ in the pole set provides at least one positive multiplicity, so the finite sum $S_\rho$ is strictly positive. For $\rho\ne0,1$ this proves $c_\rho\ne0$. For a real pole $0<\rho<1$, it also proves $c_\rho>0$. For a general nonreal pole it proves no real-order positivity. These deductions follow directly from the displayed finite sum and denominator; they do not re-prove the analytic Landau theorem.

This is a source-supported description of the final analytic mechanism: arithmetic estimates give scalar window bounds, a scalar Laplace-transform argument excludes zeros, and reflection finishes RH. It is not a claim that every upstream proof has been read. The shared factor $1/[\rho(1-\rho)]$ is exactly the scalar transform of the designated $B$ used in \fid{B-03}, since $\rho(1-\rho)=\gamma^2+1/4$. The residue calculation alone does not identify an operator representation. The further calculation in \fid{DET-11} does give an exact representation of the scalar zeta endpoint by a resolvent trace of this zero-based $C$, with all completion terms retained. The justified implication to the specified $B$ is still
\[
 \text{cited RH}\ \Longrightarrow\ \gamma_j\in\R\ \Longrightarrow\ B\ge0.
\]
Fable's reference to Landau therefore has a correct scalar counterpart in the actual endpoint. Its proposed equality of Fourier operators and Fredholm trace moments does not follow, and was explicitly refuted above. An exact historical declaration for the note's unqualified step-4 $W(1_S)$ remains unmatched; this no longer obstructs identifying the mathematical type of the examined analytic endpoint.
